%titolo{Hakuna Matata}
%autore{John, Rice}
%album{}
%tonalita{SOL}
%famiglia{Altre}
%gruppo{}
%momenti{}
%identificatore{hakuna_matata}
%data_revisione{2026_10_20}
%trascrittore{Francesco Endrici}
%video{}
\beginsong{Hakuna Matata}[by={John, Rice}]

\beginchorus
Hakuna Ma\[G]tata, ma che dolce poe\[D]sia
Hakuna Ma\[G]tata, \[E7]tutta frene\[A7]sia
Senza pen\[B-]sieri la tua vita sa\[E]rà
chi vor\[D]rà vivrà in liber\[A]tà:
Hakuna Ma\[D]tata!
\endchorus

\beginverse
Lui \[C]era una gran bel\[G]tà, \[D]
ma \[C]proprio una rari\[G]tà! \[D]
(Molto bene! Grazie!)
Ma \[F]dopo ogni pranzo lui puzzava di \[G]più,
tutti \[D]quanti svenivano e cadevano \[A]giù;
io ho l'\[C]animo sensibile, sembro \[G]tutto d'un pez\[D]zo
ma ero af\[F]franto, proprio a causa di \[G]quel mio puz\[E]zo.
E, allora che \[D]vergogna! (Si vergognava!)
E vo\[A]levo cambiar nome! (Perché soffrivi?)
Mi sen\[C]tivo sbeffeggiare! (Allora che facevi?)
(Cominciavo a scorr\dots{} No Pumba! Non vanti ai bambini. Oh scusa!)
\endverse

\beginchorus
\chordsoff
Hakuna matata, ma che dolce poesia\dots{} 
Oh, oh oh\dots{}  \rep{2}
\endchorus

\beginverse
\[fa#]Senza pen\[si-m]sieri la \[sol]tua vita sa\[mi]rà;
chi vor\[re]rà vivrà, in liber\[la]tà:
Hakuna Ma\[sol]tata\dots{}
Hakuna Ma\[re]tata.
\endverse
\endsong
